% OpenAI solução Navier--Stokes — leitura do Corpo Algébrico
% Tiffany — papers/navier-stokes.tex (secção adicionada)
\documentclass[11pt,a4paper]{article}
\usepackage[utf8]{inputenc}
\usepackage[T1]{fontenc}
\usepackage[portuguese]{babel}
\usepackage{amsmath,amssymb,amsthm}
\usepackage[margin=2.3cm]{geometry}
\usepackage{booktabs}
\usepackage[hidelinks]{hyperref}

\newcommand\code[1]{\texttt{\small #1}}
\newcommand\medido[1]{%
  \par\medskip\noindent{\small\color{gray}\textbf{Medido.} #1}\par\medskip}
\newtheorem{teorema}{Teorema}
\newtheorem{definicao}[teorema]{Definição}
\newtheorem{proposicao}[teorema]{Proposição}
\newcommand\dg{^{\dagger}}

\title{Navier--Stokes:\\
a solução OpenAI --- leitura do Corpo Algébrico}
\author{broca-so / Tiffany}
\date{Setembro 2026}
\begin{document}
\maketitle

\begin{abstract}
A OpenAI anunciou (8/Set/2026) uma prova de que as equações
Navier--Stokes tridimensionais incompressíveis podem desenvolver
uma \textbf{singularidade em tempo finito}: um vórtice que se
contrai espiralando para dentro, afinando-se como esparguete,
com velocidade ilimitada e energia finita.
A prova resolve as declarações ``C'' e ``D'' da formulação do Clay
(forçada, com força suave compacta), não as declarações ``A''/``B''
(não-forçada).
Esta secção traduz a solução para a gramática do Corpo Algébrico:
gato/esquilo, Lei~4, histerese e a fronteira renomeada.
\end{abstract}

\section{O que a OpenAI provou}\label{sec:openai}

\begin{enumerate}
\item \textbf{Configuração}: fluido em repouso inicial suave,
  força suave aplicada, viscosidade $\nu>0$ qualquer.
\item \textbf{Resultado}: existe tempo finito $T$ tal que
  $\max|u(t)|\to\infty$ quando $t\to T^{-}$,
  com energia cinética $\tfrac12\|u\|^2$ limitada durante toda a
  dinâmica — de repouso à singularidade.
\item \textbf{Mecanismo}: um vórtice que gira e se contrai;
  raio do núcleo $\sim\sqrt{T-t}$, velocidade máxima
  $\sim(T-t)^{-1/2+\varepsilon}$.
\item \textbf{Força}: o resíduo das equações é cancelado por
  pulsos espacialmente oscilatórios antes de remover os erros
  ordem a ordem — é aí que vai os 166 páginas.
\item \textbf{Formalização}: Lean (GPT-6 Astra, 17h adicionais).
\item \textbf{Escala}: $\sim$10\,000 agentes concorrentes,
  88 horas até a resolução (5/Set), mais 17h Lean.
  2,7 milhões de mensagens, 130 bilhões de tokens.
\end{enumerate}

\subsection{O que NÃO provou}
\begin{itemize}
\item Não provou as declarações ``A''/``B'' (não-forçada):
  a questão original do Clay sem força externa continua aberta.
\item Não afirma blowup com energia ilimitada — a energia fica
  finita, só a velocidade diverge.
\item A singularidade está abaixo da escala molecular antes do
  limite: o modelo continuado anuncia onde deixa de ser válido,
  não preve`` velocidade infinita``.
\end{itemize}

\section{Leitura pelo Corpo Algébrico}\label{sec:leitura}

\subsection{A equação como gato + esquilo}
Navier--Stokes incompressível:
\begin{equation}\label{eq:NS}
\partial_t u+(u\cdot\nabla)u=\nu\nabla^2 u-\nabla p,\qquad
\nabla\cdot u=0.
\end{equation}
Decompõe-se no Corpo Algébrico em:
\begin{enumerate}
\item \textbf{Stokes (linear)} = corpo parabólico = calor = \textbf{gato}:
  $A=-\nu\P\Delta$, autovalores $-\nu k^2<0$, a viscosidade dissipa.
\item \textbf{Convectivo (não-linear)} = \textbf{esquilo}:
  $\langle\P(u\cdot\nabla u),u\rangle=0$ --- conserva a energia
  (verificado $1$e$-18$); como a rotação do esquilos.
\item \textbf{Pressão/incompressibilidade} = projeção de Leray $\P$.
\end{enumerate}

\subsection{A solução como história de dois espectros}
A prova da OpenAI constrói um fluxo onde:
\begin{itemize}
\item O \textbf{esquilo} (convectivo) concentra a energia num núcleo
  que shrinks --- a energia conserva-se, só se redistribui.
\item O \textbf{gato} (viscoso) dissipa no bordo --- a viscosidade
  $\nu>0$ suaviza, mas não suficientemente depressa.
\item A \textbf{batalha} acontece na borda 3D: o esquilo ganha
  dentro do núcleo, o gato perde fora; no limite, a velocidade
  diverge e o modelo quebra.
\end{itemize}

\subsection{A força como seletor}
A força suave compacta é o \textbf{seletor} que mete o sistema
na trajetória da singularidade --- sem ela, o blowup não acontece.
No vocabulário do Corpo Algébrico: a força é a ``medida'' que
escolhe o caminho pelo projector espectral composto com o
Quantizador (Teorema~\ref{thm:maestro-memoria},
\code{tests/maestro\_memoria.js}).

\subsection{Lei~4 — a tetral}
Navier--Stokes é a Lei~4 ($T+T^{\dg}$, $|\det|=1$) lida em fluido:
o transporte incompressível é o \textbf{boost} com $|\det|=1$
(Liouville), e a regularidade em 3D é a pergunta do cruzado.
A solução da OpenAI é a metade dissipativa grampeada sobre a
metade conservativa --- exatamente a leitura que o Corpo Algébrico
nomeia como \textbf{fronteira} (docs/MAPA\_UNIVERSAL.md, §Navier--Stokes).

\subsection{Histerese e o limite contínuo}
O limite contínuo é a \textbf{histerese} (memória do caminho):
dobra+seletor sobre o tick. A solução da OpenAI constrói o
limite pelo seletor (a força) --- o caminho que leva à singularidade
é histórico, depende da força aplicada. A regularidade global 3D
é a borda aberta: o corpo aponta a fronteira, não a fecha.

\subsection{Custo mínimo e identidade de energia}
Para toda solução suave com força suave compacta:
\[
\mathcal E[u]=\int|u|^2dx
\]
decresce estritamente para $u$ não-constante. A identidade de
energia por partes em $\mathbb Z$:
\[
\sum u\,\Delta u=-\sum(u_{i+1}-u_i)^2
\]
(teorema \code{tests/navier\_corpo.js}, $8{:}0$).
A OpenAI constrói um contra-exemplo: energia finita, velocidade
ilimitada --- o gato dissipa, mas o esquilo concentra mais depressa.

\subsection{O vórtice como toro}
O vórtice esparguete é uma geometria de toro: o núcleo gira
(rotação do esquilo) e contrai (gradiente do gato). No toro
$T^2$, a acção de volume exige $|\det|=1$ (Lei~4); as ordens
do grupo de automorfismos são $\{1,2,3,4,6\}$. A contração
do núcleo é a acção de um elemento de ordem infinita --- fora
do grupo de automorfismos do toro, i.e., fora do corpo, no limite.

\section{Controvérsia e prioridade}\label{sec:prioridade}
\begin{itemize}
\item Buckmaster (NYU) e Alpöge (Anthropic) publicaram preprints
  (15/Ago) com blowup forçado para Euler 3D, equação de poros
  e Boussinesq 2D, usando LLMs (Claude, Codex/GPT-5.6 Sol, Astra).
  A OpenAI diz que não viu o trabalho antes da publicação
  (9/Set), mas Buckmaster alega que soube do método da OpenAI
  antes da execução.
\item Tao (Fields medal) chamou o resultado de Euler
  ``notável'' e perto de uma solução completa de Navier--Stokes.
\item A OpenAI não reclama o prémio do Clay.
\end{itemize}

\section{Tradução para os testes do Corpo Algébrico}\label{sec:testes}

A solução da OpenAI traduz-se nos testes existentes:
\begin{itemize}
\item \code{tests/navier\_corpo.js} ($8{:}0$): a identidade de
  energia \eqref{eq:energia} é o a-priori conhecido; a OpenAI
  mostra que é \emph{insuficiente} no 3D supercrítico.
\item \code{tests/histerese\_navier.js} ($6{:}0$): o passo de
  calor como retenção espectral --- a contração do vórtice é
  a histerese em acção (dobra+seletor).
\item \code{tests/batuta\_continuo.js} ($6{:}0$): o segmento
  afim $B(s,t)=p(t)+s\,u(t)$ --- a força é o vector director
  que leva ao blowup.
\item \code{tests/curvatura.c}: a borda $x^2=nx+1$ fecha sem
  tolerância --- o gap de Yang--Mills por andar é o comprimento
  do segmento entre ticks; na solução da OpenAI, o segmento
  encolhe a zero.
\end{itemize}

\section{Conclusão}\label{sec:conclusao}

A solução da OpenAI confirma a leitura do Corpo Algébrico:
\begin{enumerate}
\item Navier--Stokes é a \textbf{Lei~4} (tetral, $|\det|=1$).
\item A metade conservativa fecha (Liouville, medida).
\item A metade dissipativa é a \textbf{fronteira} --- o limite
  contínuo 3D é a histerese, e a regularidade global é a borda
  aberta (Clay).
\item A OpenAI constrói o contra-exemplo forçado: a singularidade
  existe \emph{com} a força, mas a questão original (sem força)
  continua aberta.
\item A gramática é uma: sete problemas, cinco operações
  (corpo\_algebrico.tex, Teorema de soluções primitivas).
\end{enumerate}

\textbf{Estaduto honesto}: a prova da OpenAI é uma claimed
resolução das declarações C/D (forçada), não das declarações
A/B (não-forçada). O Corpo Algébrico já declarava a regularidade
global 3D como fronteira aberta --- a prova da OpenAI fecha
uma das metades (a dissipativa forçada) e deixa a outra
(conservativa não-forçada) no mesmo lugar: a fronteira.

\medido{OpenAI (8/Set/2026): prova analítica + Lean formalization,
166 páginas, $\sim$10\,000 agentes, 88h + 17h Lean.
Leitura Corpo Algébrico: Lei~4, gato/esquilo, histerese,
fronteira renomeada.}
\end{document}